\section{Alexandria: A Church Between Library and Tribunal}

Origen's city was the intellectual capital of the Greek world. Roman
Alexandria kept the institutions of the Ptolemies---the scholarly community
descended from the Museum, the book trade, the grammarians and philosophers---
and added the largest Jewish diaspora community of antiquity, which had
produced in Philo a full-scale allegorical philosophy of the Septuagint a
century and a half before Origen was born. Christianity there has no recorded
founding; the later tradition of Mark's mission is not documented in the
period itself. When the Alexandrian church emerges into datable history at
the end of the second century, it already possesses the two features that
would shape Origen's whole life: learned catechesis and intermittent, lethal
persecution.

The learned catechesis is embodied in the figures Eusebius sets before
Origen: Pantaenus, a converted Stoic teacher, and Clement, the author of the
\work{Stromateis}, whom later tradition (Eusebius, and after him Photius)
made successive heads of a ``catechetical school'' in which Origen was
Clement's pupil and successor. How institutional this school really was
before Origen---an endowed office, a private philosophical circle, or simply
a succession of respected teachers---cannot now be determined, and modern
scholarship treats Eusebius's tidy succession with caution. What is certain
is that in Alexandria, uniquely, Christian instruction could look like a
school of philosophy, with a curriculum, a literature, and pagan auditors.

\witnesslimit{Nearly everything biographical in this study rests on Eusebius
of Caesarea, \work{Church History}, book 6, written about seventy to eighty
years after Origen's death by an editor of his library, co-author of a
six-book \work{Defense of Origen}, and open partisan. Eusebius worked from a
collection of more than one hundred of Origen's letters and from the reports
of ``persons still living who were acquainted with him'' (\work{Church
History} 6.2.1)---real evidence, much of it now lost---but his narrative is
a defense, not a neutral record. This study cites him constantly and trusts
him selectively, marking at each point what depends on him alone.}

The persecution came under Septimius Severus. In what Eusebius dates as ``the
tenth year of the reign of Severus''---202/203---with Laetus governing Egypt
and Demetrius newly bishop of Alexandria, ``the flame of persecution had been
kindled greatly, and multitudes had gained the crown of martyrdom''
(\work{Church History} 6.2.2--3). Among the pupils of the catechetical
circle, martyrs would soon include Plutarch, the first convert of Origen's
own teaching (\work{Church History} 6.3.2, 6.4). Whether this local storm
reflected a general imperial edict against conversion, as a later and
unreliable source asserts, is disputed; for Alexandria the fact of executions
is not. It was in this persecution that Origen's story properly begins, with
his father's arrest.
